\section{Ejercicios para el lector}

Este capítulo es un poco extenso, pero la finalidad es que usted logre identificar cuestiones básicas de un transistor así como poder comprender un circuito básico que contenga transistores.

\subsection{Preguntas teóricas}

\begin{enumerate}
  \item ¿Qué es un transistor de unión bipolar?
  \item ¿Qué tipos de transistores hay?
  \item ¿Por qué un transistor no es lo mismo que dos diodos juntos?
  \item ¿Cuántas zonas de deplexión tiene un transistor?
  \item ¿Por qué un transistor no es un capacitor cuando su base está desconectada?
  \item ¿Por qué el diseño físico y la cantidad de dopaje es importante en el diseño del transistor?
  \item ¿Qué sucede si conecto un transistor al revés?
  \item El transistor de unión bipolar ¿es un dispositivo controlado por tensión o por corriente?
  \item ¿Qué es el parámetro \(\beta\) del transistor?
  \item ¿Qué tensión es la responsable de controlar el flujo de corriente de colector?
  \item ¿Qué significa conectar un transistor en configuración emisor común?
  \item ¿Qué regiones de operación tiene un transistor?
  \item ¿Qué es la curva de carga?
  \item ¿Cómo se obtiene el punto Q?
\end{enumerate}

% Respuestas
% \begin{enumerate}
%   \item Es un dispositivo de material semiconductor conformado de tres regiones dopadas. 
%   \item Existen dos tipos: PNP y NPN
%   \item No es lo mismo que dos diodos juntos porque en un transistor no son todas sus regiones iguales. La base es muy estrecha y debilmente dopada, el colector es más grande y ligeramente dopado y el emisor es fuertemente dopado, pero más pequeño que el colector. Este diseño hace que sea posible el pasaje de cargas al colector cuando se satura la base. En un par de diodos esto nunca pasaría.
%   \item El transistor tiene dos zonas de deplexión.
%   \item El transistor no es un capacitor cuando la base está desconectada porque existen dos zonas de deplexión. Esto hace que cuando se conecte el colector al terminal positivo de una fuente de tensión y el emisor al terminal negativo, se desplacen las cargas y las zonas de deplexión, pero no se conduzca corriente eléctrica. La zona del colector se ensancha enormemente. La zona del emisor se desplaza hacia la base. Esto provoca un corrimiento de potencial de la base y consecuentemente la base tiene un potencial mayor a cero. Esto implica que no sería equivalente a un solo capacitor sino dos capacitores en serie.
%   \item El diseño físico y la cantidad de dopaje es importante por lo explicado en el inciso 3. Básicamente el transistor se diseña para que cada una de las regiones tenga un propósito específico. El emisor es un acumulador de cargas, la base es un puente o barrera según se polarice entre el emisor y colector. Y por último el colector es un gran recolector de cargas libres.
%   \item Si conecto un transistor al revés técnicamente funcionará, pero no lo hará como debería. Esto significa que su rendimiento será pobre ya que no se está aprovechando el diseño físico de los componentes.
%   \item El transistor es un dispositivo controlado por corriente.
%   \item El parámetro \(\beta\) es un parametro llamado ganancia de corriente del transistor, y se define con la relación \(I_C=\beta I_B\).
%   \item La tensión responsable de controlar \(I_C\) es \(V_{BE}\) (Base-Emisor), ya que de esta tensión depende la cantidad de corriente que circula por el colector. Con un pequeño cambio en \(V_{BE}\) se produce un cambio significativo de \(I_C\).
%   \item Significa conectar el emisor del transistor en GND o la conexión común, de modo que se formen lazos compartidos entre la entrada y la salida.
% \end{enumerate}
